\documentclass[10pt,a4paper]{amsart}
\usepackage[margin=19mm]{geometry}
\usepackage{amsmath,amssymb,mathtools}
\usepackage[scr=rsfs,cal=euler]{mathalfa}
\usepackage{xfrac}
\usepackage[unicode,hidelinks,bookmarks=false]{hyperref}
\newcommand{\ie}{i.e.\ }
\newcommand{\cf}{cf.\ }
\newcommand{\resp}{respectively\ }
\newcommand{\Subsection}{\subsection}
\providecommand{\nobreakdash}{\nobreak\hbox{-}}
\setlength{\emergencystretch}{2em}
\multlinegap0pt
\usepackage{bm}
\usepackage{bbm}
\usepackage{dsfont}
\usepackage{xspace}
\usepackage{cancel}
\usepackage{mathtools}
\usepackage{amsthm}
\makeatletter
\@ifundefined{theorem}{\newtheorem{theorem}{Theorem}}{}
\@ifundefined{lemma}{\newtheorem{lemma}[theorem]{Lemma}}{}
\@ifundefined{proposition}{\newtheorem{proposition}[theorem]{Proposition}}{}
\makeatother
\title[The structure of \'etale Boolean right restriction monoids]{The structure of \'etale Boolean right restriction monoids}
\author{Mark V. Lawson}
\date{Source: arXiv:2404.08606v3 (CC BY 4.0)}
\begin{document}
\maketitle

\noindent\emph{Source and licence.} The structure of \'etale Boolean right restriction monoids. Version arXiv:2404.08606v3 (CC BY 4.0), distributed under the Creative Commons Attribution 4.0 International licence, as recorded in the arXiv licence metadata for this exact version and shown on the abstract page https://arxiv.org/abs/2404.08606v3. Full rights notice: licenses/arxiv-2404.08606-CC-BY-4.0.txt.\par\smallskip

\noindent\textit{Editorial scope.} Counted unit: the Proposition whose source label is \texttt{prop:as} in the article (source0.tex lines 880--882, source label \texttt{prop:as}), delivered together with the article's own complete original proof of it (source0.tex lines 883--1012, one \texttt{proof} environment of 130 lines). No mathematics is written, repaired or re-derived: statement and proof are the article's own text line for line. Required context retained uncounted: lem:spiked (lines 262--270); lem:nucleus-properties (lines 855--859). Every definition, display and label of those lines is retained unchanged. Omitted from this package and not redistributed: 1--261, 271--854, 860--879, 1013--2273. Nothing inside a retained excerpt is omitted or altered: every retained body is delivered exactly as the article's lines are written, so no omission receipt of any kind accompanies this package. Citations: the retained excerpts carry no citation command at all, so no bibliography entry of the article is transcribed and no reference is redistributed. Reference adaptation: none was needed and none was made; every reference inside the delivered unit resolves inside it, so no unresolved reference or citation is produced. Printed slips kept literal: no printed word, symbol, hypothesis or number of the counted statement or of its proof was altered. Layout and class: the article's own class is \texttt{amsart} with packages including url, amsmath, cancel, graphicx, xy, eucal, amssymb, amsfonts, stackengine, xcolor; the wrapper is the lane's pinned \texttt{amsart} editorial wrapper at 10pt on A4, which restates the theorem-like environments the retained text needs and adds the article's own macro definitions that the retained text needs (none), with extra wrapper packages (bm, bbm, dsfont, xspace, cancel, mathtools). The delivered numbering is the unit's own. Assets: no retained span contains a figure environment, a graphics-inclusion command, a PDF-page inclusion, a table or a TikZ picture; no image file is copied into this package, the package declares no asset dependency and no image file is redistributed with it. Licence: version arXiv:2404.08606v3 of this article is distributed under the Creative Commons Attribution 4.0 International licence, as recorded in the arXiv licence metadata for this exact version and shown on the abstract page https://arxiv.org/abs/2404.08606v3. A full rights notice is delivered at licenses/arxiv-2404.08606-CC-BY-4.0.txt. Characters: every retained excerpt is pure ASCII, so the delivered engine, XeLaTeX with its default Latin Modern text font, reports no missing character.
\begin{lemma}\label{lem:spiked} Let $S$ be a right restriction semigroup. 
\begin{enumerate}

\item If $a,b \leq c$ and $a^{\ast} = b^{\ast}$ then $a = b$.

\item If $a,b \leq c$ then $ab^{\ast} = ba^{\ast}$.

\end{enumerate}
\end{lemma}

\begin{lemma}\label{lem:nucleus-properties} 
Let $\nu$ be a nuclus defined on the complete right restriction monoid $S$.
Then 
$$\nu (ab) = \nu (a \nu (b)) = \nu (\nu (a)b) = \nu (\nu (a) \nu (b)).$$
\end{lemma}

\begin{proposition}\label{prop:as} Let $\nu$ be a nucleus defined on the complete right restriction monoid $S$.
Then $(S_{\nu},\cdot)$ is also a complete right restriction monoid.
\end{proposition}

\begin{proof} By Lemma~\ref{lem:nucleus-properties}, $(S_{\nu},\cdot)$ is a semigroup.
It is, in fact, a monoid with identity $\nu (1)$
since
$a \cdot \nu (1) = \nu (a \nu (1))$
where $a \in S_{\nu}$.
By Lemma~\ref{lem:nucleus-properties},
we have that $\nu (a \nu (1)) = \nu (a1) = \nu (a) = a$.
We have proved that $\nu(1)$ is a right identity.
It is a left identity by symmetry.

We now prove that $S_{\nu}$ is a right restriction monoid.
Put $\mathsf{Proj} (S_{\nu}) = \{\nu (e) \colon e \in \mathsf{Proj}(S)\}$.
This is a set of projections of $S$ by axiom (N5).
Thus $\nu (1)$ is a projection.
If $a \in S_{\nu}$, define 
$$a^{\circ} = \nu (a^{\ast}).$$
This will be our `star' operation.
This is a projection by (N5).
We now show that the axioms for a right restriction semigroup hold.
Let $a,b \in S_{\nu}$.

\begin{itemize}

\item (RR1) holds:
$(a^{\circ})^{\circ} = \nu (\nu (a^{\ast})^{\ast}) = \nu (\nu (a^{\ast})) = \nu (a^{\ast}) = a^{\circ}$
by (N5) and (N3).

\item (RR2) holds: $(a^{\circ} \cdot b^{\circ})^{\circ} = (\nu(a^{\circ}b^{\circ}))^{\circ}$.
This is equal to $\nu (\nu(a^{\ast})\nu(b^{\ast}))^{\circ} = \nu (a^{\ast}b^{\ast})^{\circ}$ using Lemma~\ref{lem:nucleus-properties}.
This is equal to $\nu (\nu (a^{\ast}b^{\ast})^{\ast}) = \nu (a^{\ast}b^{\ast})$ using (N6).
Whereas, $a^{\circ} \cdot b^{\circ} = \nu (a^{\circ}b^{\circ}) = \nu (\nu (a^{\ast}) \nu (b^{\ast})) = \nu (a^{\ast}b^{\ast})$
using  Lemma~\ref{lem:nucleus-properties}.

\item (RR3) holds: $a^{\circ} \cdot b^{\circ} = \nu (a^{\circ}b^{\circ}) = \nu (b^{\circ}a^{\circ}) = b^{\circ} \cdot a^{\circ}$
where we have used (N5).

\item (RR4) holds: $a \cdot a^{\circ} = \nu (aa^{\circ}) = \nu (a\nu(a^{\ast})) = \nu (aa^{\ast}) = \nu (a) = a$ by Lemma~\ref{lem:nucleus-properties}.

\item (RR5) holds: $(a \cdot b)^{\circ} = \nu(\nu (ab)^{\ast}) = \nu ((ab)^{\ast})$ by (N6).
On the other hand, $(a^{\circ} \cdot b)^{\circ} = \nu ( \nu (\nu(a^{\ast})b)^{\ast}) = \nu (\nu (a^{\ast}b)^{\ast}) = \nu ((a^{\ast}b)^{\ast}) = \nu ((ab)^{\ast})$
by (N3) and (N6).

\item (RR6) holds:
$b^{\circ} \cdot a = \nu (b^{\circ}a) = \nu (\nu(b^{\ast})a) = \nu (b^{\ast}a)$ by Lemma~\ref{lem:nucleus-properties}.
Whereas
$a \cdot (b \cdot a)^{\circ} = \nu (a(b \cdot a)^{\circ}) = \nu (a \nu(ba)^{\circ}) = \nu (a \nu (\nu (ba)^{\ast})) = \nu (a(ba)^{\ast})$
by (N6) and  Lemma~\ref{lem:nucleus-properties}.
\end{itemize}


Thus, we have proved that $(S_{\nu},\cdot)$ is a right restriction monoid.
We now establish a number of claims.
\begin{itemize}


\item Let $a,b \in S_{\nu}$.
Denote the natural partial order on $S_{\nu}$ by $\preceq$.
Claim: $a \preceq b$ if and only if $a \leq b$.
Proof of claim.
Suppose first that $a \preceq b$.
Then $a = b \cdot a^{\circ}$.
Thus $a = \nu (ba^{\circ}) = \nu (b \nu(a^{\ast})) = \nu (ba^{\ast})$
using Lemma~\ref{lem:nucleus-properties}.
But $ba^{\ast} \leq b$ and so $\nu (ba^{\ast}) \leq \nu (b) = b$ by (N2).
Thus $a \leq b$.
Suppose now that $a \leq b$.
This means that $a = ba^{\ast}$.
Thus $a = \nu(ba^{\ast}) = \nu (b \nu (a^{\ast})) = b \cdot a^{\circ}$
by Lemma~\ref{lem:nucleus-properties}.
Whence we have proved that $a \preceq b$.

\item Claim:  $a \sim_{l} b$ in $S_{\nu}$ if and only if $a \sim_{l} b$ in $S$.
Proof of claim.
Suppose, first, that $a \sim_{l} b$ in $S_{\nu}$.
Then $a \cdot b^{\circ} = b \cdot a^{\circ}$.
This means that $\nu (a \nu (b^\ast)) = \nu (b \nu (a^{\ast}))$.
Consequently, we have that $\nu (ab^{\ast}) = \nu (ba^{\ast})$
by  Lemma~\ref{lem:nucleus-properties}.
But $ab^{\ast} \leq \nu (ab^{\ast})$ by (N1).
Similarly, $ba^{\ast} \leq  \nu (ba^{\ast})$.
It follows that $ab^{\ast} \sim_{l} ba^{\ast}$ by Lemma~\ref{lem:spiked}.
Thus $ab^{\ast}(ba^{\ast})^{\ast} = ba^{\ast}(ab^{\ast})^{\ast}$.
Whence $ab^{\ast} = ba^{\ast}$ and so $a \sim_{l} b$ in $S$.
Now, suppose that $a \sim_{l} b$ in $S$.
This means that $ab^{\ast} = ba^{\ast}$.
But $a \cdot \nu (b^{\ast}) = a \cdot b^{\circ}$
and $a \cdot \nu (b^{\ast}) = \nu (a \nu (b^{\ast})) = \nu (ab^{\ast})$ by Lemma~\ref{lem:nucleus-properties}.
It follows that $a \sim_{l} b$ in $S$.

\item Claim: if $X = \{a_{i} \colon i \in I\}$ is a left-compatible set in $S_{\nu}$
then the join in $S_{\nu}$ of $X$ exists, it is denoted by  $\bigsqcup_{i \in I} a_{i}$,
and is equal to $\nu\left(  \bigvee_{i \in I} a_{i} \right)$.
Proof of claim.
By the above, this is a left-compatible set in $S$.
It therefore has a join $a$ in $S$.
We claim that $\nu (a)$ is the join of $X$ in $S_{\nu}$.
It is an element of $S_{\nu}$ by (N2).
We have that $a_{i} \leq a$ for all $i$.
Thus  $a_{i} \leq \nu (a)$ for all $i$ by (N2).
Let $a_{i} \leq b$ for all $i$ where $b \in S_{\nu}$.
This means that $a_{i} \leq b$ in $S$.
Thus $a \leq b$.
It follows that $\nu (a) \leq b$ in $S_{\nu}$.
Thus the join of $X$ exists in $S_{\nu}$.
It follows that all joins of left-compatible subsets of $S_{\nu}$ exist.

\item Claim: $\nu \left( \bigvee_{i \in I} \nu (a_{i})  \right) = \nu (\bigvee_{i \in I} a_{i})$,
where the $a_{i}$ are arbitrary elements of $S$ which form a left-compatible set.
Proof of claim.
We have that $a_{i} \leq \bigvee_{i \in I} a_{i}$.
Thus  $\nu (a_{i}) \leq \nu \left( \bigvee_{i \in I} a_{i} \right)$ by (N2).
Thus  $\bigvee_{i \in I} \nu (a_{i}) \leq \nu \left( \bigvee_{i \in I} a_{i} \right)$.
Whence  $\nu \left(  \bigvee_{i \in I} \nu (a_{i}) \right) \leq \nu \left( \bigvee_{i \in I} a_{i} \right)$
using (N2).
To prove the reverse inequality,
we start with $a_{i} \leq \nu (a_{i})$ by (N1).
It follows that $\bigvee_{i \in I} a_{i} \leq \bigvee_{i \in I} \nu (a_{i})$ and so
$\nu \left( \bigvee_{i \in I} a_{i} \right) \leq \nu \left( \bigvee_{i \in I} \nu (a_{i}) \right)$.

\item Claim: $\left( \bigsqcup_{i \in I} a_{i} \right) \cdot b = \bigsqcup_{i \in I} a_{i} \cdot b$
in $S_{\nu}$.
Proof of claim.
We have that 
$\left( \bigsqcup_{i \in I} a_{i} \right) \cdot b = \nu \left( \bigvee_{i \in I} a_{i}b \right)$
and
$\bigsqcup_{i \in I} a_{i} \cdot b = \nu \left( \bigvee_{i \in I} \nu (a_{i}b)\right)$.
\end{itemize}

The result now follows by what we proved above.
\end{proof}

\end{document}
